We frame these results as a focused study of a single objective, with a few limitations worth stating.

\textbf{No head-to-head baselines.} We do not empirically compare against the methods we position against (feature-alignment objectives such as REPA/VIRAL~\citep{yurepresentation, yoon2025visual}, inference-time attention steering, or attention supervision in unimodal ViTs~\citep{chefer2022optimizing}); our claims concern the effect of the objective, not state-of-the-art performance.

\textbf{Intrinsic metrics measure localization, not downstream value.} AMR, AP, and NSS quantify how well attention falls on the annotated region, which is exactly what the loss optimizes (AMR especially is close to the training target); the downstream suite, not these metrics, is what tests usefulness.

\textbf{Limited scope.} Results are for one architecture (LLaVA-1.5-7B), one loss weight ($\lambda{=}0.5$), $200$ steps, and a single seed; the downstream deltas in \Cref{tab:bench} are within run-to-run noise, and a $\lambda$ sweep, longer schedules, and the squared-error variant remain untested.

\textbf{Mechanistic claims are correlational.} The drift analysis (\Cref{app:drift}) identifies candidate localization heads by how much their weights move, not by causal ablation, and reproducibility is shown across configurations rather than seeds.
